\section{Giới thiệu}
Sự bùng nổ của Internet Vạn vật (IoT) và điện toán biên đã làm gia tăng đáng kể số lượng cũng như mức độ phức tạp của các cuộc tấn công mạng. Hệ thống phát hiện xâm nhập mạng (NIDS) hiện đối mặt với ba thách thức lớn: (1) mất cân bằng dữ liệu nghiêm trọng, (2) yêu cầu phản hồi thời gian thực, và (3) hạn chế tài nguyên tại vùng biên. Các mô hình học sâu hiện tại có độ chính xác cao nhưng độ trễ suy luận lớn, trong khi các mô hình học máy truyền thống lại có tỷ lệ cảnh báo giả (FPR) cao đối với các tấn công hiếm gặp.

Để giải quyết bài toán này, chúng tôi đề xuất kiến trúc NIDS hai giai đoạn. Giai đoạn một (hướng đến độ chính xác), chúng tôi xây dựng mô hình Giáo viên (Teacher) sử dụng kiến trúc Stacking (kết hợp XGBoost, LightGBM, Random Forest), huấn luyện trên dữ liệu đã xử lý bằng chiến lược lấy mẫu lai RUS-ENN nhằm giảm thiểu tỷ lệ cảnh báo giả. Giai đoạn hai (tối ưu cho điện toán biên), kỹ thuật chưng cất tri thức (Knowledge Distillation) được áp dụng để chuyển giao tri thức từ mô hình Teacher sang một mô hình Học sinh (Student) cây quyết định nông, giúp giảm đáng kể kích thước mô hình và độ trễ suy luận.

Đóng góp của bài báo gồm: (1) một quy trình NIDS hai giai đoạn \emph{nén tri thức} của một Teacher Stacking (vốn vượt từng mô hình cơ sở một cách có ý nghĩa thống kê) vào một cây quyết định siêu nhẹ khả thi cho thiết bị biên --- giữ $\sim$99\% F1 ở mức nén 590 lần và tăng tốc 226 lần; (2) một \emph{đánh giá nghiêm ngặt có kiểm định thống kê} (McNemar, khoảng tin cậy bootstrap) so sánh chưng cất nhãn mềm với nhãn cứng và với cây huấn luyện trực tiếp, qua đó chỉ ra rằng trên CICIDS2017 một cây huấn luyện trực tiếp đã sánh ngang --- một kết quả cho thấy lợi ích của chưng cất phụ thuộc vào độ khó của dữ liệu; (3) tích hợp tiền xử lý lai RUS-ENN cùng phân tích thành phần (ablation) trên toàn bộ $\sim$2 triệu luồng của CICIDS2017.

Cấu trúc bài báo: Phần II trình bày các nghiên cứu liên quan. Phần III mô tả phương pháp đề xuất. Phần IV và V phân tích thiết lập thực nghiệm và kết quả. Phần VI thảo luận và nêu các hạn chế. Phần VII tóm tắt kết luận và hướng nghiên cứu tương lai.
